\section{实现架构、数值精度与模型边界}
\subsection{建议的 CPU/GPU 分工}
\begin{figure}[htbp]
\centering
\begin{tikzpicture}[
box/.style={draw=accent,fill=lightaccent,rounded corners,align=center,text width=40mm,minimum height=17mm,font=\small},
arr/.style={-{Latex[length=2mm]},thick,draw=accent},
node distance=12mm and 9mm]
\node[box] (cpu) {CPU：时间戳、TF、快照\\状态机与预算管理};
\node[box,right=of cpu] (maps) {GPU：版本化常驻地图\\距离场、配置空间与概率场};
\node[box,right=of maps] (batch) {GPU：候选与情景批评估\\定位、轨迹、归约};
\node[box,below=of batch] (out) {CPU：结果时效与模式接纳\\最终指令及故障处理};
\node[box,left=of out] (verify) {边界精查与实际候选复验\\独立保护路径};
\draw[arr] (cpu)--(maps);
\draw[arr] (maps)--(batch);
\draw[arr] (batch)--(out);
\draw[arr] (out)--(verify);
\end{tikzpicture}
\caption{建议的职责边界。图中只描述数据与验证关系，不声明 GPU 具有硬实时保证。最终实际执行仍须通过防护。}
\end{figure}
常驻数组优先使用 structure-of-arrays：连续存储各状态分量，
按线程访问顺序选择轨迹/时间主序。
代价图随机读取可能受带宽与缓存支配，
而不是单纯受浮点运算峰值限制。
多个极小核可融合，固定依赖图可考虑 CUDA Graphs；
但动态候选规模、重试、分支与地图更新不能为捕获方便而改变语义。
CPU 构造、内存锁和设备到主机结果回传都要计时。

输入快照必须在异步上传完成前保持有效。
不能释放 ROS 消息或修改共用数组后，仍让 GPU 读取旧指针。
所有结果带 $(\text{地图版本},\text{扫描时间},\text{姿态时间},\text{算法参数版本})$；
过期结果按已定义策略处置。
当前几何实现是同步等待流完成，不能把它描述成已经具备异步截止时限保护。

\subsection{FP32、FP64 与“更快但改变决策”的风险}
设浮点单位舍入误差为 $u$。
对 $n$ 项普通顺序求和，在 $nu<1$ 且标准模型适用时，
\[
|\operatorname{fl}(\sum_i x_i)-\sum_ix_i|
\le\gamma_{n-1}\sum_i|x_i|,\qquad
\gamma_n=\frac{nu}{1-nu}.
\]
平衡树归约的每个分量路径深度约 $\lceil\log_2n\rceil$，
可用相应深度的 $\gamma$ 估计误差；
但这只覆盖加法，不覆盖地图索引、指数和三角函数误差。
因此 CPU/GPU 不逐位相同并不必然是算法错误，
同时“总体误差小”也不能豁免边界分支检查。

建议把连续评分与硬几何分开：
\begin{itemize}
\item MPPI 轨迹与评分可先评估 FP32 的误差，并对归约采用稳定顺序/更高精度累加；
其影响用式 \eqref{eq:softbound} 及实际命令差评估。
\item AMCL 用 log-sum-exp，检查地图取格、有效束一致与后验模式稳定性。
\item 硬碰撞边界、未知区、非有限输入继续使用保守界和必要的 FP64/CPU 精查。
\item 坐标先减局部原点能减小大数相减的绝对误差，但必须在更高精度中完成该减法，
不能把两个大世界坐标先降为 FP32 再指望相减恢复丢失信息。
\end{itemize}
RTX、数据中心 GPU 与 Jetson 的 FP64 吞吐、带宽和共享内存体系不同，
本报告不把某型号的理论峰值代入实际导航加速比。
小型矩阵或标量控制没有自动使用 Tensor Core 的优势；
保持几何毫米量级语义时，未经分析的低精度矩阵化尤其不合适。

\subsection{概率场景的坐标一致与相关性}
若初始位姿和地图障碍同时受同一个全局坐标变换影响，
相对几何可保持不变。
例如共同施加刚体变换 $T$ 后
$\dist(T\mathcal F,T O)=\dist(\mathcal F,O)$。
若只在场景中扰动机器人、却把本来同样由该 TF 变换得到的障碍固定，
可能虚增相对不确定性；反之也可能漏掉地图到局部坐标的相对误差。
场景应抽样物理相对误差源，再一致地传播到机器人与障碍，
不是把所有输出协方差当独立高斯相加。

若确有局部高斯位姿误差 $e_q\sim\mathcal N(0,P_q)$，
足迹点扰动的一阶近似为
\[
e_a=J_ae_q,\quad
J_a=\begin{bmatrix}I_2&R(\theta)Ja\end{bmatrix},\quad
P_a=J_aP_qJ_a^\top.
\]
沿障碍单位法向 $n$ 的投影标准差是
$\sigma_n=\sqrt{n^\top P_an}$。
在小角度、单局部法向、高斯且协方差校准条件下，
要求名义间隙
$d_a\ge\delta+\Phi^{-1}(1-\alpha_a)\sigma_n$
可限制该一维事件的概率；$\Phi$ 为标准正态分布函数，$\alpha_a$ 是分配给该事件的概率上限。
多个时刻/障碍的联合风险可用
$\sum_a\alpha_a\le\alpha_{\rm total}$ 作保守分配，
但需有效覆盖全部所考虑事件。
该条件不适用于未校准、多峰重定位或未观测障碍；
更大的 GPU 只能更充分计算模型，不能把模型假设变成事实。

\subsection{实时性与最终保护}
CUDA 高吞吐不等于可证明的最坏响应时间。
GPU 同时被桌面、其他推理或 SLAM 使用时可能排队，
首轮上下文初始化、扩容、页面迁移都会产生尾部耗时。
预分配、预热、独立流可以降低常见开销，但不构成硬实时证书。
最终保护与失效停车应具有独立明确的输入时效规则。

若重构后的算法超期，不能无条件继续用旧控制。
需要按当前速度、当前可用障碍和已验证停止策略决定动作。
“发零速度”仍只是停车请求，不是瞬间消除动量。
报告中的所有碰撞保证均针对相应的障碍表示和运动模型，不超出原稿明确的物理边界。

\section{如何验证“算法效果提高且耗时不增加”}
\subsection{四组实验，分离硬件与算法收益}
\begin{center}\small
\begin{tabular}{p{0.12\textwidth}p{0.22\textwidth}p{0.21\textwidth}p{0.34\textwidth}}
\toprule 组别 & 算法 & 后端 & 目的\\\midrule
A & 原算法、原规模 & 优化且公平的 CPU & 冻结基线\\
B & 原算法、原规模 & CUDA & 测移植等价性与计算提速\\
C & 重构算法、扩大问题 & CPU，可离线 & 识别算法本身收益与计算代价\\
D & 同 C & CUDA，固定时延预算 & 检验用户要求的最终组合\\
\bottomrule
\end{tabular}
\end{center}
C 组离线结果只用于判断算法能力，不能拿减速仿真或暂停时钟掩盖实时超期。
A 与 D 的闭环比较必须使用相同墙钟约束、传感输入、场景难度及全部尝试分母。
先冻结候选规模选择规则，再在未用于调参的场景上比较。
若 D 使用更好的模型、更大样本、更多场景，应做逐项消融确定贡献。

\subsection{应记录的时间、规模与质量}
\begin{longtable}{p{0.19\textwidth}p{0.71\textwidth}}
\toprule 项目 & 必须记录的内容\\\midrule\endfirsthead
\toprule 项目 & 必须记录的内容\\\midrule\endhead
端到端时延 & 输入快照建立到可执行命令；P50/P95/P99、最大观测值、超期次数；首次启动与稳态分开\\
细分成本 & CPU 准备/锁、上传、核、下载、同步、地图更新、边界复核、最终候选复验\\
实际负载 & $K,H,M,N_p,B,N_o,N_q$；MPPI 调用比例；CUDA 真调用比例；每幅地图复用次数\\
时效 & 扫描、TF、地图、后验与命令年龄；是否发生丢帧、排队或旧快照复用\\
资源 & 峰值显存、GPU 时钟/功耗状态、CPU 线程数、并发其他任务、编译选项\\
定位 & 对独立真值的平移/角误差、丢定位率、错误模式率、重定位时间、后验校准与配对覆盖率\\
控制 & 全部任务成功率、完成时间、最小实物间隙、碰撞/近碰撞、跟踪误差、停止次数与指令变化率\\
风险 & 预测风险与实际发生率是否匹配；动态障碍、执行延迟、侧滑、制动能力变化的分层结果\\
\bottomrule
\end{longtable}
不将 4 mm 目标容差当作定位精度。
不将成功子集误差改善解释为整体表现改善而忽略失败分母。
不将高 GPU 利用率、kernel 时间或有限零碰撞样本替代闭环指标。

\subsection{场景与消融如何针对本报告的假设}
应覆盖：宽阔静态区、窄道与旋转角点、早停高碰撞密度、稀疏无碰撞大批量、
重复走廊多峰定位、突然全局定位偏移、动态横穿、弱纹理几何/长走廊退化、
地图更新密集、输入过期与 GPU 故障。
场景需要同时含“GPU 合适”和“CPU 早停占优”的负载，防止选择性报告。

算法消融建议：
\begin{enumerate}
\item 原几何 vs 空间索引 vs 距离场筛选 vs 配置空间，比较相同碰撞语义及保守拒绝率。
\item 同预算下只增 $K$、只增 $M$、更好动力学、结构化控制、多模式输出分别比较。
\item 同预算下更多粒子、更多束、更多独立模式、概率网格/相关评分分别比较。
\item 扩大搜索窗口、减小格距、连续精化分别比较，记录错误匹配而非只看拟合残差。
\item 固定地图更新年龄 vs 双缓冲允许年龄，量化缓存收益是否由过期地图换来。
\end{enumerate}
若同一任务重复随机运行，使用成对种子和相同输入扰动；
对 MPPI 噪声复用的基线，先固定样本以隔离后端差异，
再另做多随机种子评价。
统计不确定性应按任务/运行聚类，不能把同一条轨迹的每一时刻当独立试验。

\subsection{本次实际完成的核验}
本次未运行新导航实验或新 CUDA 基准。
报告源文件中的 \code{verify_math.py} 使用固定种子独立核验：
MPPI 熵正则驻点与权重误差界、Omni 前缀展开、
恒定旋量积分、激光/位姿图雅可比的有限差分、
距离场保守下界、方向量化位移界、
有限网格 Bayes 转移分解、FFT 相关与直接束求和、
CVaR 经验形式，以及预算/显存算例。
14 类检查全部通过，结果写入 \code{math_validation.json}。扫描雅可比有限差分最大绝对差约 $6.18\times10^{-8}$，位姿图雅可比约 $5.72\times10^{-10}$，FFT 相关与直接求和约 $7.11\times10^{-15}$；均在脚本预设容差内。运行采用独立 CPU 容器、Python 3.10 与 NumPy 1.26.4，未连接机器人或 GPU。
这些检查支持公式与数值算例的自洽，不代表新算法在机器人上可实时或更安全。
完整证明仍以文中假设与推导为准，随机数值检查不是形式化证明。

\section{建议的研究次序与可审查的成功条件}
\subsection{阶段一：拿到真实预算，并建立可比较基线}
记录整条控制路径时延与分支占比；确认实际瓶颈是 MPPI、AMCL、扫描匹配还是几何/地图更新。
先做 B 组同算法 CUDA 评估，并与带空间索引/向量化的 CPU 对照。
通过条件是：语义差异全部解释，GPU 真执行，端到端 P99 与过期输入情况可量化。
没有这一步，无法为计算扩张系数 $\alpha$ 给出可靠上限。

\subsection{阶段二：先扩大模型覆盖，再按收益分配样本}
\textbf{导航控制路线：}
在 MPPI 经常使用或尾延迟占主导时，
先实现设备常驻地图与批评分，
然后比较“增加样本”与“增加情景 + 更真实动力学 + 模式分离”。
优先增加能解释现有失败原因的维度，而不是默认把全部预算乘到 $K$ 上。
保留原直接跟踪分支作为明确基线；
若新研究要替换该分支，也单独计入调用次数与行为变化。

\textbf{定位路线：}
固定地图优先稳定 AMCL 评分与多模式保持；
歧义明显时试验粗全局/细局部分布，
只有足够复用和地图规模时再启用 FFT 全网格评分。
在线 SLAM 优先扩大有验证的相关搜索，再增加多起点连续精化，
最后才考虑大图 GPU 求解。
任何“更低匹配残差”都需独立真值检验，防止过拟合错误地图位置。

\subsection{阶段三：建立共享预计算，让更大问题变得便宜}
以静态距离场、方向配置空间和版本化动态层作为共享基础设施，
让新增航向、横移、控制情景和定位模式能够复用相同地图计算。
近边界精确查询继续保留，使粗表保守性不封锁窄道。
成功条件是\emph{包含地图构建和边界复验}后仍满足预算，而不只是查表核很快。

\subsection{最终应接受怎样的结论}
可以接受的结论是：
“在冻结版本、给定设备与规定场景分布上，
新方案相对 A 组在相同或更低端到端时延及规定超期约束下，
改善预先指定的定位/导航指标，且未恶化约定的风险和可靠性约束。”
应同时公开失败场景与未获益负载。

目前最有根据的研究主张是：
\textbf{GPU 使原来代价过高的多模式、多情景、较细配置空间搜索成为可考察方案；
真正提高效果的原因是表示、模型与搜索覆盖变得更充分，
而硬件提供了承担额外计算的可能。}
是否确实达到“效果更好且耗时不增加”，需要上述固定预算成对实验回答，
不能从显存容量、运算峰值或原稿的几何零差异测试直接推出。

\appendix
\section{原稿定位索引与复现说明}
\begin{longtable}{p{0.31\textwidth}p{0.59\textwidth}}
\toprule 原稿文件 & 本报告使用的内容\\\midrule\endfirsthead
\toprule 原稿文件 & 本报告使用的内容\\\midrule\endhead
\code{introduction.tex} & 系统模式、CUDA 现有范围与工作区版本边界\\
\code{notation.tex} & 坐标方向、底盘速度定义、填充足迹与圆盘障碍\\
\code{localization.tex} & 实际 15 维 EKF、AMCL 束乘积/KLD、Karto、图残差与质量监测\\
\code{geometry.tex} & 碰撞谓词、容差、采样界、CUDA 分块与门限\\
\code{control.tex} & 混合分支、56 步/1600 轨迹、噪声复用、原地滤波、33 条制动候选\\
\code{native_costs.tex} & 七项原生 critic、整批索引依赖、原生轮廓与填充碰撞的区别\\
\code{evidence.tex} & 历史几何差分、早停负载反向加速比、历史运动实验不能外推\\
\code{discussion.tex} & 固定间隙与概率不确定性、物理安全边界\\
\code{appendix.tex} & 有效参数、实现入口与运行平台边界\\
\code{source_manifest.json} & 本次核对的 101 个文件及 SHA-256\\
\bottomrule
\end{longtable}
本报告源目录含各节 \LaTeX{}、构建脚本、数学检查脚本、
输入与依赖来源清单以及验证结果。
使用 XeLaTeX 与 Noto CJK/URW 字体可重建 PDF；
数值检查需要 Python 3 与 NumPy。
原始压缩包未被改写，项目源文件和配置未被本任务修改。
